%% JAIR submission. Preamble, title block, abstract, acknowledgements and the
%% reproducibility checklist follow the JAIR Author Kit (JAIR_Example_Template.tex,
%% 2025-09-15). The body comes from paper/*.md via paper/jair/build.py.
%% Style commands are unchanged from the kit, as JAIR requires.
\documentclass[manuscript, screen, review]{jair}

\setcopyright{cc}
\acmDOI{10.1613/jair.1.xxxxx}

\JAIRAE{Insert JAIR AE Name}
\JAIRTrack{}
\acmVolume{4}
\acmArticle{6}
\acmMonth{8}
\acmYear{2026}

\RequirePackage[
  datamodel=acmdatamodel,
  style=acmauthoryear,
  backend=biber,
  giveninits=true,
  uniquename=init
  ]{biblatex}

\renewcommand*{\bibopenbracket}{(}
\renewcommand*{\bibclosebracket}{)}

\addbibresource{references.bib}

%% Packages the pandoc-generated tables need.
\usepackage{longtable}
\usepackage{calc}
\usepackage{array}
\providecommand{\tightlist}{\setlength{\itemsep}{0pt}\setlength{\parskip}{0pt}}
\graphicspath{{../../results/}}

\begin{document}

\title[Model, Not Mode]{Model, Not Mode: A Pre-registered Model x Mode Grid on a Constrained Design Task}

\author{Roshan Aryal}
\authornote{Corresponding Author.}
\email{roshanaryaal@gmail.com}
\affiliation{%
  \institution{Independent Researcher}
  \city{Auckland}
  \country{New Zealand}
}

\renewcommand{\shortauthors}{Aryal}

\begin{abstract}
{\bf Background:}
Chat apps let users choose a model and also how hard it thinks. Leaderboards name the model
but rarely the mode, so it is unclear which choice matters more on open-ended work.

{\bf Objectives:}
We test whether the reasoning mode explains more of the variance in answer quality than the
model family does, on a demanding design task with hard constraints.

{\bf Methods:}
Six model families were crossed with two modes, fast and reasoning, with five answers per
cell collected in each vendor's own chat app (60 answers). Three blind raters, one human and
two LLMs, scored every answer on a nine-dimension rubric with web verification. A
pre-registered two-way decomposition compared the variance explained by mode and by family,
with bootstrap intervals and permutation tests.

{\bf Results:}
Family explained far more variance in the total score than mode (difference in omega
squared -0.527, 95\% CI [-0.666, -0.417]), so the confirmatory hypothesis was not
supported. Reasoning scored 1.37 points higher out of 18 (95\% CI [0.77, 1.93]), but most of
the gain came from one family and it halved once answer length was removed. Rater agreement
fell below the pre-registered bar on seven of nine dimensions, so the secondary hypotheses
are exploratory.

{\bf Conclusions:}
On this task, choosing the model matters more than choosing how hard it thinks.
Comparisons should still report the mode, because its effect differs between models.
\end{abstract}

\maketitle

% Generated by paper/latex_build.py from paper/*.md. Do not edit by hand.

\section{Introduction}

Leaderboards rank models. People who use these models choose more than a model. In every
major chat app they also choose how hard the model thinks: Instant or High in ChatGPT,
Effort Low or High in Claude, Fast or Expert in Grok, DeepThink on or off in DeepSeek. The
same weights can then behave like a different system. If the mode switch changes answer
quality as much as switching vendor does, a comparison that names the model but not the
mode is confounded, and a user who picks a vendor carefully but leaves the mode on its
default may be optimising the smaller factor.

Prior work on reasoning modes mostly measures accuracy on tasks with an answer key, such as
math, code and science questions, and usually within one model family or one training
recipe (Section 2). Real requests are often open-ended and constrained: there is no single
right answer, but there are hard limits a good answer must respect. For these tasks we do
not know how much of the variation in quality comes from the mode and how much from the
model.

We test this directly. We hold one demanding prompt fixed: design a fully local, always-on
autonomous AI workstation for one stated machine, a Mac mini with 32 GB of unified memory.
We cross six model families with two modes, fast and reasoning, and collect five answers
per cell in each vendor's own web app, for 60 answers. Three blind raters, one human and
two LLMs from different vendors, score every answer on a nine-dimension rubric with
mechanical checks for the memory budget and internal consistency and web verification for
every named tool and model. A pre-registered two-way variance decomposition then asks, per
dimension, whether mode explains more of the variance than family does, and whether
reasoning scores higher.

The question came from an observation. LA-P1 \cite{aryal2026lap1} captured one model family in
three modes and saw answers that looked different in ways that mattered, but with one answer
per mode and no controls its adjudicated scores could not separate a mode effect from
run-to-run noise, and did not show a reasoning advantage. This paper is the controlled test.

Our contributions are:

\begin{enumerate}
\def\labelenumi{\arabic{enumi}.}
\tightlist
\item
  A crossed 6 x 2 model-by-mode design with five replicates per cell, run in the chat apps
  people use, with an API check that each mode switch changes how much the model reasons.
\item
  A pre-registered analysis (OSF, frozen before the first capture) that decomposes
  answer-quality variance into family, mode, interaction and residual per rubric
  dimension, using bias-corrected omega squared, bootstrap intervals and permutation tests,
  with one confirmatory hypothesis.
\item
  Evidence on rubric portability across a grid the rubric's authors never saw:
  inter-rater reliability per dimension and agreement between raters on how the 12 cells
  rank, which a follow-up benchmark (LA-P5) needs.
\item
  An open harness, \texttt{modegrid}, that covers planning, capture, blinding, scoring and
  analysis, and the full data: 60 verbatim answers, blinded packets' build rules, all rater
  sheets and every departure from the plan.
\end{enumerate}

Family, not mode, dominated. On the total score, family explained far more variance than
mode (delta omega squared -0.527, 95\% CI {[}-0.666, -0.417{]}), so the confirmatory hypothesis
is not supported. Reasoning still scored higher than fast, by 1.37 points of 18 (95\% CI
{[}0.77, 1.93{]}), but the gain came mostly from one family, Qwen, and halved once answer
length was removed. Family dominated eight of nine rubric dimensions, and the ninth was
inconclusive. Per-answer agreement between raters was low on most dimensions, yet the
three raters ranked the 12 cells almost identically on the total (Kendall's W 0.955). For
evaluation practice: name the model first, but report the mode too, because what it does
differs from one model to the next.

\section{Related work}

\subsection{Thinking modes in hybrid models}

Hybrid models expose reasoning as a switch on one set of weights. Work on these models asks
how to switch and whether the switch is clean. HRBench compares prompt-based, routing and
speculative switching strategies on six models and five math, science and code benchmarks
\cite{ning2026hrbench}. Several studies show that the ``off'' mode is not fully off. \citet{wangs2025hybrid} find only partial separation between think and no-think modes, with reasoning
behaviour leaking into no-think answers. \citet{liu2026leakage} attribute 42\% to 79\% of
no-think gains to leaked thinking behaviour, and Path-Lock Expert separates the two modes
at the architecture level to stop it \cite{wangs2026pathlock}. Thinking persists even under
explicit instructions not to think, more so as the answer space opens \cite{lei2026inertia}. These results bear on our design in two ways. A fast mode may still reason on an
open-ended task, so we measure reasoning tokens for each mode through the API (Section
3.4). And three of our six families offer only a lowest setting, not an off switch, which
a pre-registered robustness check separates out.

The closest design to ours compares instant and thinking modes on the same checkpoints
across five frontier vendors, for moral judgments \cite{madur2026moral}. It finds that verdict agreement barely changes with thinking. It also
notes, as we do, that some vendors' lowest setting is not a clean off. That study works on
verdicts with a small answer space, through the API, and
measures agreement between models without separate raters. Ours works on an
open-ended design task with hard constraints, in the chat apps people use, with a human
rater among three blind raters.

\subsection{Model choice versus other factors}

A growing line of work asks how much of the variation in LLM output comes from the model
and how much from something around it. On creative tasks, a variance decomposition across
12 models and 10 prompts finds that model choice explains about as much variance in
originality as the prompt, and more in fluency \cite{haase2026creative}. For coding agents, the
harness that runs the model can explain more variance than the model and can reverse model
rankings \cite{zhang2026harness}. Test-time compute work shows that extra inference compute,
spent well, can beat a larger model on math \cite{snell2025scaling}. We apply the same
variance-decomposition logic to a factor none of these studies varies: the reasoning mode
of a fixed model, crossed with model identity, on a constrained task.

\subsection{Reasoning and factuality}

Reasoning does not only help. Reasoning-oriented RL fine-tuning raises hallucination rates
\cite{li2025hallucinate}, reasoning LLMs hallucinate more than non-reasoning models on long-form
factuality \cite{chen2026factuality}, and thinking can turn a correct direct answer into a wrong
one \cite{wangk2026margo}. On closed-book knowledge tasks, more test-time computation
does not reliably raise accuracy and often adds hallucinations \cite{zhao2026tts}. Reasoning-model errors on constrained problems often come
from inventing problem features that were never given \cite{heyman2025hallucinating}. These
works compare differently trained models or study training methods, not one model with
reasoning switched on and off at inference time. They motivate a competing prediction for
our hypothesis H3: if the effect carries over, reasoning mode would lower tool factuality
(D3) rather than leave it unchanged.

\subsection{Run-to-run variance and LLM raters}

Repeated runs of one model with fixed settings give different outputs. Sources include
hosted-service non-determinism at temperature 0 \cite{atil2025nondeterminism}, floating-point
non-associativity that shifts reasoning-model accuracy by up to 9\% \cite{yuan2025numerical}, and
early divergence in agent trajectories that moves single-run scores by 2 to 6 points
\cite{bjarnason2026randomness}. Across 12 models, 10 reasoning strategies and 30 trials
each, about three-quarters of score variance comes from benchmark, system and item
structure, with a residual that single runs absorb \cite{potamitis2025reasonbench}. We therefore collect five replicates per cell and treat
replicate variance as the residual. LLM judges score the same input differently across
models and settings \cite{lau2026judge}, and closed-book raters flag real tools released after
their training cutoff as fabricated \cite{aryal2026lap2}. We use three blind raters, one of them human,
take the median, require web verification for tool and model factuality, and report each
LLM rater's scores for its own vendor's family separately.

\subsection{LLMs as systems architects}

LLMs are increasingly asked to design systems. On networked systems, \citet{sahu2026architects}
find that LLMs cannot yet do architectural reasoning reliably and that how the LLM is
involved changes the outcome a lot. LA-P1 \cite{aryal2026lap1} introduced the task and rubric we
reuse, and captured one family (DeepSeek) in three modes, one answer each, with no
controls. Its adjudicated totals did not show a reasoning advantage (reasoning 8, fast 9,
of 18). That pilot cannot separate a mode effect from run-to-run noise. This paper is the
controlled test.

\subsection{What is new here}

To our knowledge, no prior study crosses reasoning mode with model identity and decomposes
answer-quality variance on an open-ended, constrained design task. We add a pre-registered
6 x 2 design with replicates, a per-dimension decomposition with bias-corrected effect
sizes, blinded scoring by human and LLM raters with web verification, and a reusable
harness. The search was last repeated on 2026-10-26.

\section{Method}

\subsection{Design overview}

We ran a crossed 6 x 2 design: six model families (Claude, GPT, Gemini, Qwen, DeepSeek,
Grok) by two modes (fast, reasoning), with five independent replicates per cell, for 60
scored answers. Every answer responds to the same fixed prompt. Within a family, the model
is held fixed and only the vendor's reasoning control changes. The study was
pre-registered before the first scored capture: the plan was frozen on 2026-10-13 (git tag
\texttt{prereg-v1}) and registered on OSF the same day (\url{https://osf.io/nj2zp/}). The single
confirmatory test (H4) asks whether mode explains more variance in the total rubric score
than family does, and whether reasoning scores higher than fast. Three secondary tests
(H1 to H3) ask the same of single rubric dimensions (section 3.6).

\subsection{Task and prompt}

The task is the frozen prompt \texttt{prompt-v1} from LA-P1 \cite{aryal2026lap1}, reused unchanged and
pinned by SHA-256 (\texttt{ba2acba8...0370e0}). It asks the model to act as a senior AI
infrastructure architect and design a fully local, always-on autonomous AI workstation for
one stated machine: an Apple Mac mini with a base M6 chip, 32 GB unified memory, a 512 GB
internal SSD and a 1 TB external SSD. The prompt fixes the date as 31 August 2026, asks for
research into ten areas (local inference, agent runtimes, multi-agent architecture, coding
agents, research agents, memory, computer control, always-on operation, remote access,
storage), and requires eleven labelled deliverables (A to K), including a memory plan (C),
an installation plan (I) and a list of what not to install (J). It is 1,782 words long.

The task has no answer key. It has hard constraints that can be checked: the memory
budget, internal consistency between the recommendations and deliverable J, and whether
each named tool and model exists. That makes it a test of grounding and constraint
reasoning rather than recall. The scope of every conclusion is this one task.

\subsection{Models and modes}

For each vendor we chose a current general-purpose chat model, available to the author's
individual account in the vendor's own web app, that offers both a fast and a reasoning
setting. Table 1 lists the settings. The choice was fixed in the pre-registration before any
capture. These are not all each vendor's top model, and they are not matched on size or
price.

\textbf{Table 1.} Models and mode settings. Chat settings were used for the 60 scored runs. API
settings were used once per cell to check that the switch works (section 3.4). The
pre-registration listed GPT-5.6 Sol and Grok 4.6. Both apps had moved to newer models by the
time of collection, GPT-6 and Grok 4.7, and both ran the same model in both modes (Appendix
D).

{\def\LTcaptype{none} % do not increment counter
\begin{longtable}[]{@{}
  >{\raggedright\arraybackslash}p{(\linewidth - 14\tabcolsep) * \real{0.1250}}
  >{\raggedright\arraybackslash}p{(\linewidth - 14\tabcolsep) * \real{0.1250}}
  >{\raggedright\arraybackslash}p{(\linewidth - 14\tabcolsep) * \real{0.1250}}
  >{\raggedright\arraybackslash}p{(\linewidth - 14\tabcolsep) * \real{0.1250}}
  >{\raggedright\arraybackslash}p{(\linewidth - 14\tabcolsep) * \real{0.1250}}
  >{\raggedright\arraybackslash}p{(\linewidth - 14\tabcolsep) * \real{0.1250}}
  >{\raggedright\arraybackslash}p{(\linewidth - 14\tabcolsep) * \real{0.1250}}
  >{\raggedright\arraybackslash}p{(\linewidth - 14\tabcolsep) * \real{0.1250}}@{}}
\toprule\noalign{}
\begin{minipage}[b]{\linewidth}\raggedright
Family
\end{minipage} & \begin{minipage}[b]{\linewidth}\raggedright
Model
\end{minipage} & \begin{minipage}[b]{\linewidth}\raggedright
Web app
\end{minipage} & \begin{minipage}[b]{\linewidth}\raggedright
Fast (chat)
\end{minipage} & \begin{minipage}[b]{\linewidth}\raggedright
Reasoning (chat)
\end{minipage} & \begin{minipage}[b]{\linewidth}\raggedright
API fast
\end{minipage} & \begin{minipage}[b]{\linewidth}\raggedright
API reasoning
\end{minipage} & \begin{minipage}[b]{\linewidth}\raggedright
Fast floor
\end{minipage} \\
\midrule\noalign{}
\endhead
\bottomrule\noalign{}
\endlastfoot
Claude & Claude Sonnet 5.5 & claude.ai & Effort Low & Effort High & adaptive thinking, effort low & adaptive thinking, effort high & lowest available \\
GPT & GPT-6 & chatgpt.com (paid) & Thinking effort Instant & Thinking effort High & \texttt{reasoning\_effort:\ none} & \texttt{reasoning\_effort:\ high} & off \\
Gemini & Gemini 3.8 Flash & aistudio.\allowbreak{}google.com & Thinking level Low & Thinking level High & \texttt{reasoning\_effort:\ low} & \texttt{reasoning\_effort:\ high} & lowest available \\
Qwen & Qwen3.8-Max & chat.qwen.ai & Fast & Thinking & \texttt{enable\_thinking:\ false} & \texttt{enable\_thinking:\ true} & off \\
DeepSeek & DeepSeek-V4.1-Flash & chat.\allowbreak{}deepseek.com & DeepThink off & DeepThink on & \texttt{enable\_thinking:\ false} & \texttt{enable\_thinking:\ true} & off \\
Grok & Grok 4.7 & grok.com & Fast & Expert & \texttt{reasoning\_effort:\ low} & \texttt{reasoning\_effort:\ xhigh} & lowest available \\
\end{longtable}
}

\textbf{Two kinds of contrast.} Three families can switch reasoning fully off (GPT, Qwen,
DeepSeek). The other three can only set it to a minimum (Claude Sonnet 5.5 rejects
\texttt{thinking:\ disabled}, Gemini 3.8 Flash has no level below low, Grok 4.7 cannot disable
reasoning). For those three the contrast is ``low versus high reasoning'', not ``off versus
on''. A pre-registered robustness check analyses the two groups separately (section 3.7).

\textbf{Same model in both modes.} We excluded any switch that changes the model (Flash versus
Pro, mini versus full) and any automatic router (``Auto'' in ChatGPT, Grok, Qwen and the
consumer Gemini app). Gemini was therefore run in Google AI Studio, Google's own web app,
where the model is picked by name. The consumer Gemini app offered only ``Auto'', which can
route between Flash-Lite, Flash and Pro by itself. For Claude, GPT, Gemini and Qwen the app
showed the same model name in both modes. The author confirmed that grok.com ran Grok 4.7 in
both Fast and Expert. chat.\allowbreak{}deepseek.com shows no model name, and the DeepSeek model is the
author's observation of what the app served on 2026-10-13. For DeepSeek we therefore describe
the contrast as the vendor's exposed chat setting, not a verified same-weights manipulation,
as the pre-registration requires.

\textbf{Model self-reports.} The prompt asks each model to name itself. Every GPT answer named
``GPT-6'' and every Grok answer ``Grok 4.7'', which is correct. Every DeepSeek and Gemini answer
named an Anthropic Claude model, which is wrong: the app's model picker (Gemini) and the
author's observation (DeepSeek) are the evidence of which model ran. Every self-report is
kept in its capture.

\subsection{Checking each mode switch through the API}

Chat apps do not report reasoning tokens. To show that each documented switch changes how
much the model reasons, we made one API call per family and mode (12 calls) with the
documented API equivalent of each chat setting (Table 1) and the same prompt. These calls
are evidence about the switch only: they were never scored, blinded or pooled with the
chat answers, and they cannot show what a chat app runs internally.

\textbf{Table 2.} API mode check. Reasoning tokens as reported by each API (or derived where the
API omits them).

{\def\LTcaptype{none} % do not increment counter
\begin{longtable}[]{@{}
  >{\raggedright\arraybackslash}p{(\linewidth - 8\tabcolsep) * \real{0.2000}}
  >{\raggedright\arraybackslash}p{(\linewidth - 8\tabcolsep) * \real{0.2000}}
  >{\raggedright\arraybackslash}p{(\linewidth - 8\tabcolsep) * \real{0.2000}}
  >{\raggedright\arraybackslash}p{(\linewidth - 8\tabcolsep) * \real{0.2000}}
  >{\raggedright\arraybackslash}p{(\linewidth - 8\tabcolsep) * \real{0.2000}}@{}}
\toprule\noalign{}
\begin{minipage}[b]{\linewidth}\raggedright
Family
\end{minipage} & \begin{minipage}[b]{\linewidth}\raggedright
Model returned
\end{minipage} & \begin{minipage}[b]{\linewidth}\raggedright
Reasoning tokens, fast
\end{minipage} & \begin{minipage}[b]{\linewidth}\raggedright
Reasoning tokens, reasoning
\end{minipage} & \begin{minipage}[b]{\linewidth}\raggedright
Verdict
\end{minipage} \\
\midrule\noalign{}
\endhead
\bottomrule\noalign{}
\endlastfoot
Claude & \texttt{claude-sonnet-5-5} & 0 & 14,864 & switch confirmed \\
GPT & \texttt{gpt-5.6-sol} & 0 & 6,000 & switch confirmed \\
Gemini & \texttt{gemini-3.8-flash} & 0 & 4,911 & switch confirmed \\
Qwen & \texttt{qwen3.8-max} & 0 & 23,805 & switch confirmed \\
DeepSeek & \texttt{deepseek-v4.1-flash} & 0 & 1,367 & switch confirmed \\
Grok & \texttt{grok-4.6} & 131 & 4,627 & lowest available \\
\end{longtable}
}

The DeepSeek and Qwen checks ran on Alibaba Model Studio, a third-party host. The GPT
reasoning call spent its whole 6,000-token output limit (set by the account's rate limit)
on reasoning and returned no visible answer, and the Claude reasoning call stopped at its
32,000-token limit. Neither affects the evidence that the switch works. The GPT and Grok
checks ran on the pre-registered models (\texttt{gpt-5.6-sol}, \texttt{grok-4.6}), not on the GPT-6 and
Grok 4.7 that the apps served, so for those two families they show the documented switch on
the planned model only. A second check on 2026-10-19 used the served models. GPT-6 (\texttt{gpt-6-sol}) returned 0
reasoning tokens in fast and 6,000 in reasoning, so the switch is confirmed on the model
the app ran. Grok 4.7 (\texttt{grok-4.7}) answered in full at effort low with 1,648 reasoning
tokens, but at effort xhigh it stopped twice after a one-sentence preamble, with about
300 to 360 reasoning tokens, as if waiting for a search tool the API call did not offer.
We made no third attempt and changed no setting, so the second Grok check is
inconclusive, and Grok stays in the lowest-available group on the first check. Claude and Gemini
show zero reasoning tokens at their lowest setting on this call, although their
documentation says reasoning cannot be disabled, so we keep them in the ``lowest available''
group.

\subsection{Collecting the answers}

All 60 scored answers were collected in each vendor's web app, the way people use these
models and the way LA-P1 was collected. The run order came from a manifest randomised with
seed 20261001. Families were captured one block at a time, in the order each first appears
in the manifest (Qwen, Grok, DeepSeek, Claude, Gemini, GPT). Inside a block, runs followed
manifest order, which mixes fast and reasoning runs, so a vendor change during a block does
not line up with one mode. Each block had to finish within 48 hours and the whole
collection within 14 days.

Each run used a new chat, signed in, with memory off, no custom instructions and no files.
The prompt text was pasted unchanged. Web search was turned off where the app allows it
(DeepSeek's Search toggle and AI Studio's Google Search grounding, both on by default).
Three apps offer no way to turn search off in the chosen mode, and their models searched the
web by themselves: Claude, GPT and Grok. We record this as app behaviour (section 3.8). A
capture file was created immediately before each send with \texttt{modegrid\ new-capture}, which
stamps the UTC time in a field that is never edited. Where a send time was not recorded
exactly (Qwen runs 3 to 10, sent before their capture files were made, and the Gemini and
GPT blocks, whose chats ran side by side), the capture gives the send time as an interval
bounded by known events instead of a guess. Each capture
stores the final answer verbatim, the model name the app showed, the interface settings,
and, where the app shows one, the reasoning panel. Raters never see the reasoning panel.

The Qwen, Gemini and GPT blocks were sent by the author by hand. The Grok, DeepSeek and
Claude blocks were sent, at the author's request, by an AI assistant operating the author's
signed-in browser through a browser extension, with the same settings and one new chat per
run (Appendix D). Collection took place from 2026-10-13 10:45 UTC to
2026-10-14 13:08 UTC, about 26 hours. Every block finished well inside 48 hours: Qwen took
12.9 hours, Claude 4.6, Grok 1.8, GPT 1.1, Gemini 0.6 and DeepSeek 0.2.

\textbf{Replacements.} A run was replaced only for an app error, an answer cut off in
transmission, or the wrong model or mode. It was never replaced for a low score or an odd
answer. The replacement was captured straight away in the same block and kept the run id.
The failed attempt was kept in a separate folder and never scored. Getting 60 valid answers
took 66 attempts: one DeepSeek answer cut off at the app's output limit, one DeepSeek app
error, one Qwen run sent in the wrong mode by the operator, and three Claude reasoning runs
stopped by the account's five-hour usage limit. The pre-registration lists errors, cut-offs
and wrong models as reasons to replace. We treated an operator choosing the wrong mode the
same way as a wrong model, and we treated the usage limit as the app being unavailable, so
those three redos did not count toward the two-replacement limit per cell. One extra Qwen
fast chat was sent by mistake and is kept but assigned to no run. No family was dropped.
No answer was an outright refusal: no consensus score is 0 on all nine dimensions. Every
departure from the frozen plan is listed with its date and reason in Appendix D.

\subsection{Scoring}

\textbf{Rubric.} Answers were scored on the nine-dimension rubric from LA-P1, each dimension
0, 1 or 2 (total 0 to 18): D1 hardware-constraint adherence, D2 recency, D3 tool
factuality, D4 model factuality, D5 benchmark factuality, D6 citation quality, D7
actionability, D8 security model, and D9 internal consistency. We used the v2 hard anchors
for D1, D7, D8 and D9, which were written to fix low agreement on those four dimensions in
LA-P1. Two checks are mechanical. For D1 the rater runs the memory checker
\texttt{memory\_budget.py} from LA-P1 on the set of models the answer keeps loaded at the same
time. For D9 the rater first compares deliverable J (``what not to install'') with the rest
of the answer, and recommending something J forbids scores 0. An outright refusal scores 0
on every dimension.

\textbf{Web verification.} For D3 and D4 the rater searches the web for every recommended tool,
framework and model before scoring. A name that resolves is real. A name that cannot be
found either way is unresolved and scores 1, not 0. Only a name that a search shows does not
exist, or a wrong load-bearing attribute (such as model size), scores 0. We added this rule
because closed-book raters in LA-P1 flagged real tools released after their training cutoff
as fabricated \cite{aryal2026lap2}.

\textbf{Raters.} The author collected the captures, so the author does not score. Three raters
score every answer independently and blind: one human who is not an author, and two LLMs
from different vendors, each with live web search: \texttt{llm-a}, Claude Opus 5.5 (Anthropic),
and \texttt{llm-b}, GPT-6 at Thinking effort High (OpenAI). The human rater (Rater H2) is not an
author and asked that neither their name nor their background be published; the author
knows their identity. Rater H2 replaced an earlier human rater whose returns
were made with AI tools and were not used (Appendix D). \texttt{llm-b} scored on 2026-10-19,
\texttt{llm-a} from 2026-10-19 to 2026-10-20, and Rater H2 from 2026-10-24 to 2026-10-26. Both LLM raters come from vendors whose
families are scored, so each one's scores for its own vendor's family are also reported
separately.

\textbf{Sessions.} Each rater scored one answer per session, from a single file that holds the
fixed rater prompt, the rubric, the scoring protocol, the memory checker and one blinded
answer. Each LLM rater used a new chat for every answer with memory, history and personal
connectors off and web search on, and returned one CSV row. Each rater had their own
random reading order. All three raters got the same instructions. The rubric, reused
unchanged from LA-P1, still carries notes about that study: one names a model, and another
describes what earlier raters missed in certain answers. A rule in the rater prompt told raters
to ignore those parts.

\textbf{Blinding.} Raters saw only the final answer. Reasoning panels, interface labels,
timestamps and capture metadata were removed. Self-identification lines near the top of an
answer (the model's own name, which the prompt asks for) were removed without a trace,
and a vendor or model name was masked only where the answer refers to itself. Names used as content were
kept, because the prompt asks for a model and tool stack and D2 to D4 cannot be scored with
the recommendations hidden. A conflict-of-interest disclosure in which an answer names its
own maker (``so discount my view on'' that vendor's product) was removed as well. We also
removed clutter that the chat apps add to copied text and that would name the family
outright, such as ChatGPT's citation chips and Grok's citation tags. The original captures
are never edited. Before scoring, an automated audit
(\texttt{modegrid\ blinding-audit}) checked every packet for app clutter, self-disclosure and
masked self-references, and found none. It also counted the model's own style habits that
survive blinding: four habits each appear in only one family (tool-progress text and notes
on the real date in Claude answers, ``as of my cutoff'' in DeepSeek answers, LaTeX arrows in
Gemini answers), and each is split evenly between fast and reasoning runs, so it can hint
at family but not at mode. Raters guessed the family and mode of every answer, and we
report how often they beat chance as the blinding-failure rate (section 3.7).

\textbf{Consensus.} The analysis unit is the median of the three raters' scores for each answer
and dimension. There is no hand adjudication, because the only person who could adjudicate,
the author, is not blind. The total for an answer is the sum of its nine median scores.

\textbf{Reliability.} We report Krippendorff's ordinal alpha per dimension across the three
raters, and interval alpha on each rater's total. The pre-registered bar is 0.667. A
dimension below the bar is reported in full but treated as exploratory. H4 stays
confirmatory if the total score's alpha clears the bar. A supported result with alpha
between 0.667 and 0.800 is described as tentative.

\subsection{Analysis}

\textbf{Primary statistic.} For each dimension and for the total, we fit a balanced two-way
fixed-effects decomposition (family, mode, family x mode, residual) on the 60 consensus
scores, and compute

delta = omega2(mode) - omega2(family),

where omega2(effect) = (SS\_effect - df\_effect x MS\_residual) / (SS\_total + MS\_residual).
Negative estimates are kept. We use omega squared rather than eta squared because family
has 5 degrees of freedom and mode has 1, and eta squared rewards extra degrees of freedom
even when they carry only noise. The direction of the mode effect is the mean of reasoning
minus the mean of fast.

\textbf{Uncertainty.} We use a 95\% percentile bootstrap with 5,000 resamples (seed 20261001),
resampling the five replicates within each cell. Families are treated as fixed, so
conclusions are about these six families. Permutation p-values (5,000 permutations) are
reported beside each interval. The mode permutation shuffles scores between modes within
each family, so it tests the strict null that mode makes no difference in any family. Where
the interval and the p-value disagree, we report both. The bootstrap interval decides every
verdict.

\textbf{Decision rule.} ``Mode \textgreater{} model'' on a dimension means delta is positive and its 95\% CI
excludes zero. ``Model \textgreater{} mode'' means delta is negative and its CI excludes zero. Anything
else is inconclusive. H1, H2 and H4 also require reasoning to score higher than fast, with
a 95\% CI above zero. If 90\% or more of a dimension's scores share one value, the dimension
is reported as uninformative and its hypothesis as inconclusive.

\textbf{Hypotheses.}

\begin{itemize}
\tightlist
\item
  H4 (confirmatory): on the total score, mode \textgreater{} model and reasoning \textgreater{} fast.
\item
  H1 (secondary): on D1 hardware-constraint adherence, mode \textgreater{} model and reasoning \textgreater{} fast.
\item
  H2 (secondary): on D9 internal consistency, mode \textgreater{} model and reasoning \textgreater{} fast.
\item
  H3 (secondary): on D3 tool factuality, model \textgreater{} mode.
\end{itemize}

H4 is the only confirmatory test and is not corrected for multiplicity. H1 to H3 are
reported uncorrected, with that caveat. The other six dimensions are descriptive. For every
dimension we also report omega squared for the family x mode interaction and each family's
reasoning-minus-fast difference. If the interaction explains more than both main effects,
the H4 verdict is labelled ``average effect only''.

\textbf{Robustness checks} (pre-registered, reported whatever the result):

\begin{enumerate}
\def\labelenumi{\arabic{enumi}.}
\tightlist
\item
  The same analysis on rank-transformed scores.
\item
  The same analysis after removing the linear effect of log answer length in words. Mode
  changes length, so a shrinking mode effect would mean the two cannot be separated, not
  that length causes the effect.
\item
  The same analysis on each rater's own scores.
\item
  The same analysis on the three true-off families (GPT, Qwen, DeepSeek) and on the three
  lowest-available families (Claude, Gemini, Grok). With 30 runs each, these are
  descriptive.
\end{enumerate}

\textbf{Blinding failure.} For each rater, we test whether family and mode guesses beat chance
with an exact one-sided binomial test (mode against 1/2, family against 1/6). If any rater
beats chance on mode, H1 to H4 are re-run on each rater's scores and we report whether the
raters who did not beat chance reach the same conclusions. If all three beat chance, a
supported H4 is labelled tentative.

\textbf{Power.} Simulation (\texttt{modegrid\ power}, Appendix) gives power of about 0.9 to detect a
mode effect whose variance is five times the family variance, and about 0.4 to 0.6 for a
moderate effect. Under equal effects, 15\% of simulated studies still show ``mode \textgreater{} model'',
which we report beside the H4 result.

\subsection{Departures from the plan that bear on interpretation}

The Deviations table (Appendix D) lists every change. Four bear on how to read the results:

\begin{itemize}
\tightlist
\item
  \textbf{Web search.} Claude, GPT and Grok searched the web by themselves in every run, in both
  modes, because their apps offer no switch for it. Gemini and DeepSeek ran with search
  off, and Qwen with no search option turned on. Differences between families can therefore
  include differences in live search, which is how these apps are used.
\item
  \textbf{Operators.} Three blocks were sent by the author and three by an AI assistant in the
  author's browser. Settings and procedure were the same.
\item
  \textbf{Model identity.} The GPT and Grok apps ran newer models than pre-registered, GPT-6
  and Grok 4.7, released shortly before collection. Each ran one model in both modes, so
  both contrasts remain mode switches, but their API checks cover the planned models.
  DeepSeek shows no model name, so its contrast is the vendor's exposed setting.
\item
  \textbf{Blinding fixes before scoring.} Four leaks were found and fixed before any real
  scoring: app clutter copied with the answers, bare ``Model:'' label lines missed by the
  self-identification rule, the marker lines left where header lines were removed, whose
  count and position differed by family, and conflict-of-interest disclosures that named
  the answer's maker through its own product even with the name masked. All four fixes
  apply to every packet in the same way, and the blinding key did not change.
\end{itemize}

\subsection{Materials and code}

The harness (\texttt{modegrid}), configuration, captures, packets' build rules, scores and analysis
outputs are released at {[}GitHub tag and Zenodo DOI{]}. The prompt, rubric and memory checker
are vendored from LA-P1 at commit \texttt{d5e39e7} and pinned by SHA-256. All randomness (run
order, reading orders, bootstrap, permutations) is seeded, so the analysis reproduces
exactly from the released scores with \texttt{modegrid\ analyze\ -\/-seed\ 20261001\ -\/-boot\ 5000\ -\/-perm\ 5000}.

\section{Results}

\subsection{Data and reliability}

All 60 planned answers were scored by all three raters (the human rater, \texttt{llm-a} and
\texttt{llm-b}), and all 60 have a consensus score. No answer has a consensus of 0 on all nine
dimensions (\texttt{n\_all\_zero} is 0 in every cell of \texttt{results/cell\_means.csv}), so no answer was
treated as an outright refusal.

\textbf{Table 3.} Inter-rater reliability, three raters: Krippendorff's alpha, ordinal per
dimension and interval on the total.

{\def\LTcaptype{none} % do not increment counter
\begin{longtable}[]{@{}lll@{}}
\toprule\noalign{}
Dimension & alpha & Status \\
\midrule\noalign{}
\endhead
\bottomrule\noalign{}
\endlastfoot
D1 hardware-constraint & 0.574 & below 0.667 \\
D2 recency & 0.283 & below 0.667 \\
D3 tool-factuality & 0.249 & below 0.667 \\
D4 model-factuality & 0.630 & below 0.667 \\
D5 benchmark-factuality & 0.599 & below 0.667 \\
D6 citation-quality & 0.539 & below 0.667 \\
D7 actionability & 0.858 & ok \\
D8 security-model & 0.862 & ok \\
D9 internal-consistency & 0.285 & below 0.667 \\
Total & 0.698 & ok \\
\end{longtable}
}

Seven of nine dimensions fall below the 0.667 bar: D1, D2, D3, D4, D5, D6 and D9. Their
tests are exploratory. This includes all three secondary hypotheses (H1 on D1, H2 on D9,
H3 on D3). Only D7 and D8 clear the bar. The total clears 0.667 (alpha 0.698), so H4 stays
confirmatory. It does not clear 0.800, so any supported result on the total would be
tentative.

Pairwise, the human rater and \texttt{llm-a} have a mean exact agreement of 0.72 over D1 to D9
(range 0.18 to 0.93 across dimensions). The human rater and \texttt{llm-b} have 0.71 (range 0.45
to 0.95).

\subsection{Confirmatory test (H4)}

On the total score (Table 4, ``total'' row):

\begin{itemize}
\tightlist
\item
  delta omega squared (mode minus family) is -0.527, 95\% CI {[}-0.666, -0.417{]};
\item
  omega squared is 0.050 for mode and 0.577 for family;
\item
  reasoning scores 1.37 points higher than fast out of 18, 95\% CI {[}0.77, 1.93{]};
\item
  permutation p is 0.007 for mode and below 0.001 for family. Both agree with their
  intervals;
\item
  omega squared for the family x mode interaction is 0.184. It exceeds the mode effect
  (0.050) but not the family effect (0.577), so it does not exceed both main effects and
  the ``average effect only'' label does not apply.
\end{itemize}

Verdict: \textbf{model \textgreater{} mode}. H4 status: \textbf{no}. H4 predicted mode \textgreater{} model with reasoning
higher than fast. Reasoning did score higher than fast, with a CI above zero, but family
explained far more of the variance in the total than mode did, and the CI on delta lies
wholly below zero.

Power caveat: under equal mode and family effects, 15\% of simulated studies still show
``mode \textgreater{} model'' (Section 3.7). The observed result is in the opposite direction, so this
caveat does not weaken it, but it is reported as pre-registered.

\subsection{Secondary tests (H1 to H3) and the other dimensions}

\textbf{Table 4.} Variance decomposition, all nine dimensions and the total
(\texttt{results/variance\_table.md}). Delta is omega squared for mode minus omega squared for
family, with a 95\% bootstrap CI over replicates.

{\def\LTcaptype{none} % do not increment counter
\begin{longtable}[]{@{}
  >{\raggedright\arraybackslash}p{(\linewidth - 24\tabcolsep) * \real{0.0769}}
  >{\raggedright\arraybackslash}p{(\linewidth - 24\tabcolsep) * \real{0.0769}}
  >{\raggedright\arraybackslash}p{(\linewidth - 24\tabcolsep) * \real{0.0769}}
  >{\raggedright\arraybackslash}p{(\linewidth - 24\tabcolsep) * \real{0.0769}}
  >{\raggedright\arraybackslash}p{(\linewidth - 24\tabcolsep) * \real{0.0769}}
  >{\raggedright\arraybackslash}p{(\linewidth - 24\tabcolsep) * \real{0.0769}}
  >{\raggedright\arraybackslash}p{(\linewidth - 24\tabcolsep) * \real{0.0769}}
  >{\raggedright\arraybackslash}p{(\linewidth - 24\tabcolsep) * \real{0.0769}}
  >{\raggedright\arraybackslash}p{(\linewidth - 24\tabcolsep) * \real{0.0769}}
  >{\raggedright\arraybackslash}p{(\linewidth - 24\tabcolsep) * \real{0.0769}}
  >{\raggedright\arraybackslash}p{(\linewidth - 24\tabcolsep) * \real{0.0769}}
  >{\raggedright\arraybackslash}p{(\linewidth - 24\tabcolsep) * \real{0.0769}}
  >{\raggedright\arraybackslash}p{(\linewidth - 24\tabcolsep) * \real{0.0769}}@{}}
\toprule\noalign{}
\begin{minipage}[b]{\linewidth}\raggedright
Dim
\end{minipage} & \begin{minipage}[b]{\linewidth}\raggedright
eta2 family
\end{minipage} & \begin{minipage}[b]{\linewidth}\raggedright
eta2 mode
\end{minipage} & \begin{minipage}[b]{\linewidth}\raggedright
eta2 FxM
\end{minipage} & \begin{minipage}[b]{\linewidth}\raggedright
eta2 resid
\end{minipage} & \begin{minipage}[b]{\linewidth}\raggedright
omega2 FxM
\end{minipage} & \begin{minipage}[b]{\linewidth}\raggedright
delta omega2
\end{minipage} & \begin{minipage}[b]{\linewidth}\raggedright
95\% CI
\end{minipage} & \begin{minipage}[b]{\linewidth}\raggedright
reasoning - fast
\end{minipage} & \begin{minipage}[b]{\linewidth}\raggedright
95\% CI
\end{minipage} & \begin{minipage}[b]{\linewidth}\raggedright
p mode
\end{minipage} & \begin{minipage}[b]{\linewidth}\raggedright
p family
\end{minipage} & \begin{minipage}[b]{\linewidth}\raggedright
Verdict
\end{minipage} \\
\midrule\noalign{}
\endhead
\bottomrule\noalign{}
\endlastfoot
D1 & 0.206 & 0.000 & 0.147 & 0.647 & 0.079 & -0.150 & {[}-0.386, -0.045{]} & 0.000 & {[}-0.200, 0.200{]} & 1.000 & 0.035 & model \textgreater{} mode \\
D2 & 0.692 & 0.004 & 0.063 & 0.241 & 0.037 & -0.664 & {[}-0.840, -0.547{]} & 0.067 & {[}-0.067, 0.200{]} & 0.662 & 0.000 & model \textgreater{} mode \\
D3 & 0.524 & 0.011 & 0.084 & 0.381 & 0.044 & -0.477 & {[}-0.710, -0.321{]} & -0.100 & {[}-0.233, 0.033{]} & 0.450 & 0.000 & model \textgreater{} mode \\
D4 & 0.415 & 0.029 & 0.161 & 0.395 & 0.119 & -0.350 & {[}-0.637, -0.175{]} & 0.167 & {[}0.000, 0.300{]} & 0.187 & 0.000 & model \textgreater{} mode \\
D5 & 0.582 & 0.039 & 0.125 & 0.254 & 0.098 & -0.519 & {[}-0.656, -0.440{]} & 0.333 & {[}0.133, 0.533{]} & 0.041 & 0.000 & model \textgreater{} mode \\
D6 & 0.654 & 0.000 & 0.013 & 0.333 & -0.022 & -0.622 & {[}-0.805, -0.494{]} & 0.000 & {[}-0.133, 0.133{]} & 1.000 & 0.000 & model \textgreater{} mode \\
D7 & 0.187 & 0.058 & 0.213 & 0.542 & 0.155 & -0.083 & {[}-0.300, 0.008{]} & 0.200 & {[}0.033, 0.333{]} & 0.096 & 0.045 & inconclusive \\
D8 & 0.403 & 0.074 & 0.134 & 0.389 & 0.093 & -0.294 & {[}-0.485, -0.186{]} & 0.267 & {[}0.100, 0.400{]} & 0.020 & 0.000 & model \textgreater{} mode \\
D9 & 0.316 & 0.066 & 0.226 & 0.392 & 0.184 & -0.216 & {[}-0.409, -0.095{]} & 0.433 & {[}0.167, 0.700{]} & 0.034 & 0.000 & model \textgreater{} mode \\
total & 0.595 & 0.053 & 0.200 & 0.152 & 0.184 & -0.527 & {[}-0.666, -0.417{]} & 1.367 & {[}0.767, 1.933{]} & 0.007 & 0.000 & model \textgreater{} mode \\
\end{longtable}
}

p = 0.000 means p \textless{} 0.001 at 5,000 permutations.

All three secondary tests are exploratory, because their dimensions fall below the alpha
bar (Table 3). They are reported uncorrected for multiplicity, as pre-registered, and
cannot by themselves carry a claim.

\begin{itemize}
\tightlist
\item
  \textbf{H1} (D1, predicted mode \textgreater{} model and reasoning \textgreater{} fast): model \textgreater{} mode, delta -0.150
  {[}-0.386, -0.045{]}. Reasoning minus fast is 0.00 {[}-0.20, 0.20{]}. The prediction fails in
  direction. Exploratory (alpha 0.57).
\item
  \textbf{H2} (D9, predicted mode \textgreater{} model and reasoning \textgreater{} fast): model \textgreater{} mode, delta -0.216
  {[}-0.409, -0.095{]}. Reasoning scores higher than fast, 0.43 {[}0.17, 0.70{]}, but family still
  explains more variance than mode. The prediction fails in direction. Exploratory
  (alpha 0.29).
\item
  \textbf{H3} (D3, predicted model \textgreater{} mode): model \textgreater{} mode, delta -0.477 {[}-0.710, -0.321{]}.
  Reasoning minus fast is -0.10 {[}-0.23, 0.03{]}. Of the pre-stated outcomes, model \textgreater{} mode
  occurred, which supports H3. The competing hallucination account (mode \textgreater{} model with
  reasoning lower) did not occur. Exploratory (alpha 0.25), so the support is weak.
\end{itemize}

For every hypothesis the permutation p for mode agrees with its interval.

Other dimensions, descriptive: D2, D4, D5, D6 and D8 are model \textgreater{} mode. D7 is inconclusive
(delta -0.083 {[}-0.300, 0.008{]}). No dimension had no variance or was uninformative in the
full design. Delta is negative on every dimension.

The random-effects variance components are in Appendix B, labelled exploratory.

\subsection{Where mode mattered: per-family differences}

Figure~\ref{fig:cells} shows cell means, fast vs reasoning, one line per family,
for the total, D1, D9 and D3.

\begin{figure}[t]
  \centering
  \includegraphics[width=\linewidth]{figure1.pdf}
  \caption{Cell means, fast versus reasoning, one line per family, for the total score and
  for D1, D9 and D3. Consensus scores, five answers per cell.}
  \Description{Line plots of mean scores in fast and reasoning mode for six model families.}
  \label{fig:cells}
\end{figure}

\textbf{Table 5.} Total score (out of 18) by family and mode, consensus means over five
replicates (\texttt{results/family\_mode\_effects.csv}).

{\def\LTcaptype{none} % do not increment counter
\begin{longtable}[]{@{}llll@{}}
\toprule\noalign{}
Family & Fast & Reasoning & Reasoning - fast \\
\midrule\noalign{}
\endhead
\bottomrule\noalign{}
\endlastfoot
Claude & 13.80 & 15.60 & +1.80 \\
DeepSeek & 11.80 & 11.20 & -0.60 \\
Gemini & 9.40 & 9.40 & 0.00 \\
GPT & 15.00 & 16.40 & +1.40 \\
Grok & 13.60 & 12.40 & -1.20 \\
Qwen & 6.80 & 13.60 & +6.80 \\
\end{longtable}
}

Three of six families (Claude, GPT, Qwen) moved in the same direction as the average
(+1.37). Gemini did not move, and DeepSeek and Grok scored lower in reasoning. Qwen's +6.80
is the largest single difference; without Qwen the mean of the other five is +0.28. The
interaction did not exceed both main effects (Section 4.2), so the pre-registration does
not require interpreting these differences. They are reported as descriptive.

\subsection{Robustness}

\textbf{Table 6.} Robustness checks 1 and 2: delta omega squared on raw scores and on ranks, and
delta, omega squared for mode and reasoning minus fast after removing the linear effect of
log answer length (\texttt{results/variance\_table.md}, sensitivity checks).

{\def\LTcaptype{none} % do not increment counter
\begin{longtable}[]{@{}
  >{\raggedright\arraybackslash}p{(\linewidth - 16\tabcolsep) * \real{0.1111}}
  >{\raggedright\arraybackslash}p{(\linewidth - 16\tabcolsep) * \real{0.1111}}
  >{\raggedright\arraybackslash}p{(\linewidth - 16\tabcolsep) * \real{0.1111}}
  >{\raggedright\arraybackslash}p{(\linewidth - 16\tabcolsep) * \real{0.1111}}
  >{\raggedright\arraybackslash}p{(\linewidth - 16\tabcolsep) * \real{0.1111}}
  >{\raggedright\arraybackslash}p{(\linewidth - 16\tabcolsep) * \real{0.1111}}
  >{\raggedright\arraybackslash}p{(\linewidth - 16\tabcolsep) * \real{0.1111}}
  >{\raggedright\arraybackslash}p{(\linewidth - 16\tabcolsep) * \real{0.1111}}
  >{\raggedright\arraybackslash}p{(\linewidth - 16\tabcolsep) * \real{0.1111}}@{}}
\toprule\noalign{}
\begin{minipage}[b]{\linewidth}\raggedright
Dim
\end{minipage} & \begin{minipage}[b]{\linewidth}\raggedright
delta (raw)
\end{minipage} & \begin{minipage}[b]{\linewidth}\raggedright
delta (ranks)
\end{minipage} & \begin{minipage}[b]{\linewidth}\raggedright
Same sign
\end{minipage} & \begin{minipage}[b]{\linewidth}\raggedright
delta (length-adj.)
\end{minipage} & \begin{minipage}[b]{\linewidth}\raggedright
omega2 mode (raw)
\end{minipage} & \begin{minipage}[b]{\linewidth}\raggedright
omega2 mode (length-adj.)
\end{minipage} & \begin{minipage}[b]{\linewidth}\raggedright
Reasoning - fast (raw)
\end{minipage} & \begin{minipage}[b]{\linewidth}\raggedright
Reasoning - fast (length-adj.)
\end{minipage} \\
\midrule\noalign{}
\endhead
\bottomrule\noalign{}
\endlastfoot
D1 & -0.150 & -0.145 & yes & -0.120 & -0.013 & -0.004 & 0.000 & -0.093 \\
D2 & -0.664 & -0.664 & yes & -0.618 & -0.001 & 0.012 & 0.067 & 0.128 \\
D3 & -0.477 & -0.477 & yes & -0.432 & 0.003 & -0.008 & -0.100 & 0.011 \\
D4 & -0.350 & -0.341 & yes & -0.346 & 0.021 & 0.022 & 0.167 & 0.171 \\
D5 & -0.519 & -0.462 & yes & -0.575 & 0.034 & 0.013 & 0.333 & 0.224 \\
D6 & -0.622 & -0.645 & yes & -0.575 & -0.007 & 0.018 & 0.000 & -0.153 \\
D7 & -0.083 & -0.079 & yes & -0.124 & 0.046 & -0.014 & 0.200 & 0.001 \\
D8 & -0.294 & -0.294 & yes & -0.180 & 0.065 & -0.002 & 0.267 & -0.080 \\
D9 & -0.216 & -0.219 & yes & -0.232 & 0.057 & 0.051 & 0.433 & 0.412 \\
total & -0.527 & -0.570 & yes & -0.655 & 0.050 & 0.008 & 1.367 & 0.621 \\
\end{longtable}
}

\begin{enumerate}
\def\labelenumi{\arabic{enumi}.}
\item
  \textbf{Ranks.} Delta on ranks has the same sign as on raw scores on every dimension and the
  total. The total moves from -0.527 to -0.570. The verdict does not change.
\item
  \textbf{Length.} After removing log answer length, the mode effect on the total falls by more
  than half: reasoning minus fast drops from 1.37 to 0.62 points, and omega squared for mode
  from 0.050 to 0.008. Model \textgreater{} mode still holds (delta -0.655). As pre-registered, we
  report that the mode effect is not separable from answer length, not that length causes
  it. Reasoning answers are longer, and longer answers score higher; this design cannot
  tell the two apart. The effect also shrinks on D3, D5, D7 and D8, and almost vanishes on
  D7 and D8. It does not shrink on D9 (0.43 to 0.41), the one dimension where the gain from
  reasoning survives the adjustment. The first version of this check judged shrinking on
  delta rather than on the mode effect; the Deviations table records the correction
  (Appendix D).
\item
  \textbf{Per rater.} Each rater's own scores give the same H4 verdict as the consensus.

  {\def\LTcaptype{none} % do not increment counter
  \begin{longtable}[]{@{}
    >{\raggedright\arraybackslash}p{(\linewidth - 8\tabcolsep) * \real{0.2000}}
    >{\raggedright\arraybackslash}p{(\linewidth - 8\tabcolsep) * \real{0.2000}}
    >{\raggedright\arraybackslash}p{(\linewidth - 8\tabcolsep) * \real{0.2000}}
    >{\raggedright\arraybackslash}p{(\linewidth - 8\tabcolsep) * \real{0.2000}}
    >{\raggedright\arraybackslash}p{(\linewidth - 8\tabcolsep) * \real{0.2000}}@{}}
  \toprule\noalign{}
  \begin{minipage}[b]{\linewidth}\raggedright
  Rater
  \end{minipage} & \begin{minipage}[b]{\linewidth}\raggedright
  H1 (D1)
  \end{minipage} & \begin{minipage}[b]{\linewidth}\raggedright
  H2 (D9)
  \end{minipage} & \begin{minipage}[b]{\linewidth}\raggedright
  H3 (D3)
  \end{minipage} & \begin{minipage}[b]{\linewidth}\raggedright
  H4 (total)
  \end{minipage} \\
  \midrule\noalign{}
  \endhead
  \bottomrule\noalign{}
  \endlastfoot
  Human & model \textgreater{} mode, -0.10 {[}-0.30, 0.10{]} & inconclusive, 0.37 {[}0.13, 0.60{]} & model \textgreater{} mode, 0.17 {[}0.10, 0.23{]} & model \textgreater{} mode, 1.57 {[}1.03, 2.10{]} \\
  \texttt{llm-a} & model \textgreater{} mode, 0.00 {[}-0.20, 0.20{]} & model \textgreater{} mode, 0.30 {[}0.10, 0.53{]} & model \textgreater{} mode, -0.13 {[}-0.30, 0.03{]} & model \textgreater{} mode, 1.23 {[}0.63, 1.80{]} \\
  \texttt{llm-b} & model \textgreater{} mode, 0.07 {[}-0.10, 0.23{]} & model \textgreater{} mode, 0.20 {[}-0.13, 0.53{]} & model \textgreater{} mode, -0.07 {[}-0.23, 0.10{]} & model \textgreater{} mode, 1.20 {[}0.53, 1.87{]} \\
  \end{longtable}
  }

  Cells give the verdict, then reasoning minus fast with its 95\% CI. All three raters give
  model \textgreater{} mode on H4, with reasoning above fast and a CI above zero. On H2 the human
  rater's verdict is inconclusive where the other two give model \textgreater{} mode. On H3 the human
  rater scored reasoning higher than fast on D3 (0.17 {[}0.10, 0.23{]}), where the two LLM
  raters show no clear difference.
\item
  \textbf{Fast-mode floor.} The same analysis on the true-off families (GPT, Qwen, DeepSeek)
  and the lowest-available families (Claude, Gemini, Grok), 30 runs each, descriptive.

  {\def\LTcaptype{none} % do not increment counter
  \begin{longtable}[]{@{}
    >{\raggedright\arraybackslash}p{(\linewidth - 8\tabcolsep) * \real{0.2000}}
    >{\raggedright\arraybackslash}p{(\linewidth - 8\tabcolsep) * \real{0.2000}}
    >{\raggedright\arraybackslash}p{(\linewidth - 8\tabcolsep) * \real{0.2000}}
    >{\raggedright\arraybackslash}p{(\linewidth - 8\tabcolsep) * \real{0.2000}}
    >{\raggedright\arraybackslash}p{(\linewidth - 8\tabcolsep) * \real{0.2000}}@{}}
  \toprule\noalign{}
  \begin{minipage}[b]{\linewidth}\raggedright
  Group
  \end{minipage} & \begin{minipage}[b]{\linewidth}\raggedright
  Measure
  \end{minipage} & \begin{minipage}[b]{\linewidth}\raggedright
  delta (95\% CI)
  \end{minipage} & \begin{minipage}[b]{\linewidth}\raggedright
  Reasoning - fast (95\% CI)
  \end{minipage} & \begin{minipage}[b]{\linewidth}\raggedright
  Verdict
  \end{minipage} \\
  \midrule\noalign{}
  \endhead
  \bottomrule\noalign{}
  \endlastfoot
  True-off & total & -0.348 {[}-0.554, -0.161{]} & 2.53 {[}1.67, 3.40{]} & model \textgreater{} mode \\
  True-off & D1 & -0.167 {[}-0.425, 0.018{]} & -0.07 {[}-0.40, 0.27{]} & inconclusive \\
  True-off & D9 & 0.004 {[}-0.291, 0.283{]} & 0.73 {[}0.40, 1.07{]} & inconclusive \\
  Lowest-available & total & -0.736 {[}-0.853, -0.617{]} & 0.20 {[}-0.60, 1.00{]} & model \textgreater{} mode \\
  Lowest-available & D1 & -0.163 {[}-0.538, 0.066{]} & 0.07 {[}-0.13, 0.27{]} & inconclusive \\
  Lowest-available & D9 & -0.357 {[}-0.612, -0.156{]} & 0.13 {[}-0.27, 0.53{]} & model \textgreater{} mode \\
  \end{longtable}
  }

  H1, H2 and H4 hold in neither group, so the pre-registered true-off rule is not
  triggered. Both groups give model \textgreater{} mode on the total. The mode effect on the total is
  larger where fast mode turns reasoning off (+2.53, CI above zero) than where it only
  lowers it (+0.20, CI includes zero), and delta is closer to zero in the true-off group.
  The reasoning-over-fast gain in the full design comes mainly from the true-off families,
  Qwen above all (Table 5).
\end{enumerate}

\subsection{Blinding}

\textbf{Table 7.} Blinding check: share of answers guessed, guess accuracy against chance, and
the exact one-sided binomial p-value (\texttt{results/blinding\_check.csv}).

{\def\LTcaptype{none} % do not increment counter
\begin{longtable}[]{@{}llllll@{}}
\toprule\noalign{}
Rater & Guess & Guessed & Accuracy & Chance & p \\
\midrule\noalign{}
\endhead
\bottomrule\noalign{}
\endlastfoot
Human & family & 60 of 60 & 0.15 & 0.17 & 0.688 \\
Human & mode & 60 of 60 & 0.70 & 0.50 & 0.001 \\
\texttt{llm-a} & family & 60 of 60 & 0.75 & 0.17 & \textless{} 0.001 \\
\texttt{llm-a} & mode & 60 of 60 & 0.68 & 0.50 & 0.003 \\
\texttt{llm-b} & family & 60 of 60 & 0.20 & 0.17 & 0.292 \\
\texttt{llm-b} & mode & 60 of 60 & 0.60 & 0.50 & 0.077 \\
\end{longtable}
}

Two raters beat chance on mode: the human rater (0.70, p = 0.001) and \texttt{llm-a} (0.68,
p = 0.003). \texttt{llm-b} did not (0.60, p = 0.077). As pre-registered, H1 to H4 were re-run on
each rater's scores (Section 4.5, check 3). \texttt{llm-b}, the rater who did not beat chance,
reaches the same H4 verdict as the consensus (model \textgreater{} mode). Not all three raters beat
chance, so the ``mode blinding failed'' label does not apply. H4 is not supported in any
case.

On family, \texttt{llm-a} guessed right on 75\% of answers (chance 17\%, p \textless{} 0.001). The human rater
(15\%) and \texttt{llm-b} (20\%) did not beat chance. Family leaking through recommended model names
is expected (Appendix D2) and changes no hypothesis's status. It does mean that \texttt{llm-a}
correctly identified the family of most answers it scored. Section 4.8 reports how each
LLM rater scored its own vendor's answers.

One leak was found after scoring. Three packets, all Claude reasoning answers, still
contained the author's account name or home city, which the claude.ai app had passed to
the model although memory was off; the pre-scoring audit had no probe for it. All three
raters guessed ``reasoning'' for all three packets, and \texttt{llm-a} guessed Claude for all three
(the human rater and \texttt{llm-b} for one each). Three of 60 answers cannot move the rates in
Table 7 by more than 0.05; no scores were changed (Appendix D).

\subsection{Rubric portability (exploratory, for LA-P5)}

\textbf{Table 8.} Agreement of the three raters on the ranking of the 12 cell means: Kendall's W
across raters and Kendall's tau-b for each pair (\texttt{results/cell\_rank\_agreement.csv}).

{\def\LTcaptype{none} % do not increment counter
\begin{longtable}[]{@{}
  >{\raggedright\arraybackslash}p{(\linewidth - 8\tabcolsep) * \real{0.2000}}
  >{\raggedright\arraybackslash}p{(\linewidth - 8\tabcolsep) * \real{0.2000}}
  >{\raggedright\arraybackslash}p{(\linewidth - 8\tabcolsep) * \real{0.2000}}
  >{\raggedright\arraybackslash}p{(\linewidth - 8\tabcolsep) * \real{0.2000}}
  >{\raggedright\arraybackslash}p{(\linewidth - 8\tabcolsep) * \real{0.2000}}@{}}
\toprule\noalign{}
\begin{minipage}[b]{\linewidth}\raggedright
Measure
\end{minipage} & \begin{minipage}[b]{\linewidth}\raggedright
Kendall's W
\end{minipage} & \begin{minipage}[b]{\linewidth}\raggedright
Human \textasciitilde{} \texttt{llm-a}
\end{minipage} & \begin{minipage}[b]{\linewidth}\raggedright
Human \textasciitilde{} \texttt{llm-b}
\end{minipage} & \begin{minipage}[b]{\linewidth}\raggedright
\texttt{llm-a} \textasciitilde{} \texttt{llm-b}
\end{minipage} \\
\midrule\noalign{}
\endhead
\bottomrule\noalign{}
\endlastfoot
D1 & 0.891 & 0.87 & 0.67 & 0.75 \\
D2 & 0.675 & 0.72 & 0.39 & 0.19 \\
D3 & 0.544 & -0.05 & 0.10 & 0.61 \\
D4 & 0.685 & 0.48 & 0.34 & 0.63 \\
D5 & 0.842 & 0.51 & 0.65 & 0.80 \\
D6 & 0.875 & 0.93 & 0.62 & 0.68 \\
D7 & 0.948 & 0.83 & 0.90 & 0.94 \\
D8 & 0.951 & 0.86 & 0.79 & 0.95 \\
D9 & 0.870 & 0.63 & 0.61 & 0.85 \\
total & 0.955 & 0.88 & 0.81 & 0.80 \\
\end{longtable}
}

The raters rank the 12 cells almost the same way on the total (W = 0.955), even though
their per-answer agreement on several dimensions is low (Table 3). D3 is the weakest: the
human rater's ranking is unrelated to either LLM rater's (tau-b -0.05 and 0.10), while the
two LLM raters agree moderately (0.61). D2 and D4 are also low. No dimension had no
variance, so none is unusable as a benchmark dimension on that ground.

\subsection{Own-vendor scores}

\textbf{Table 9.} Each LLM rater on its own vendor's family (\texttt{results/own\_vendor.csv}). Own gap
is the rater's mean score minus the mean of the other two raters on its own vendor's
answers. Baseline gap is the same on the other five families. Self-preference is own gap
minus baseline gap, with a 95\% CI. Descriptive: no hypothesis rests on it.

{\def\LTcaptype{none} % do not increment counter
\begin{longtable}[]{@{}
  >{\raggedright\arraybackslash}p{(\linewidth - 10\tabcolsep) * \real{0.1667}}
  >{\raggedright\arraybackslash}p{(\linewidth - 10\tabcolsep) * \real{0.1667}}
  >{\raggedright\arraybackslash}p{(\linewidth - 10\tabcolsep) * \real{0.1667}}
  >{\raggedright\arraybackslash}p{(\linewidth - 10\tabcolsep) * \real{0.1667}}
  >{\raggedright\arraybackslash}p{(\linewidth - 10\tabcolsep) * \real{0.1667}}
  >{\raggedright\arraybackslash}p{(\linewidth - 10\tabcolsep) * \real{0.1667}}@{}}
\toprule\noalign{}
\begin{minipage}[b]{\linewidth}\raggedright
Rater
\end{minipage} & \begin{minipage}[b]{\linewidth}\raggedright
Family
\end{minipage} & \begin{minipage}[b]{\linewidth}\raggedright
Measure
\end{minipage} & \begin{minipage}[b]{\linewidth}\raggedright
Own gap
\end{minipage} & \begin{minipage}[b]{\linewidth}\raggedright
Baseline gap
\end{minipage} & \begin{minipage}[b]{\linewidth}\raggedright
Self-preference (95\% CI)
\end{minipage} \\
\midrule\noalign{}
\endhead
\bottomrule\noalign{}
\endlastfoot
\texttt{llm-a} & Claude & total & 0.30 & -0.81 & 1.11 {[}0.52, 1.71{]} \\
\texttt{llm-a} & Claude & D1 & -0.05 & 0.00 & -0.05 {[}-0.18, 0.06{]} \\
\texttt{llm-a} & Claude & D2 & 0.35 & -0.09 & 0.44 {[}0.25, 0.62{]} \\
\texttt{llm-a} & Claude & D3 & 0.05 & -0.34 & 0.39 {[}0.10, 0.66{]} \\
\texttt{llm-a} & Claude & D4 & 0.00 & -0.02 & 0.02 {[}-0.13, 0.18{]} \\
\texttt{llm-a} & Claude & D5 & 0.05 & 0.26 & -0.21 {[}-0.63, 0.18{]} \\
\texttt{llm-a} & Claude & D6 & 0.25 & 0.00 & 0.25 {[}0.01, 0.47{]} \\
\texttt{llm-a} & Claude & D7 & 0.05 & 0.04 & 0.01 {[}-0.19, 0.26{]} \\
\texttt{llm-a} & Claude & D8 & 0.00 & -0.03 & 0.03 {[}-0.03, 0.09{]} \\
\texttt{llm-a} & Claude & D9 & -0.40 & -0.63 & 0.23 {[}-0.18, 0.67{]} \\
\texttt{llm-b} & GPT & total & 0.10 & -1.40 & 1.50 {[}0.70, 2.34{]} \\
\texttt{llm-b} & GPT & D1 & 0.10 & -0.15 & 0.25 {[}0.01, 0.51{]} \\
\texttt{llm-b} & GPT & D2 & 0.05 & -0.33 & 0.38 {[}0.05, 0.69{]} \\
\texttt{llm-b} & GPT & D3 & 0.15 & -0.12 & 0.27 {[}0.06, 0.51{]} \\
\texttt{llm-b} & GPT & D4 & -0.10 & -0.12 & 0.02 {[}-0.23, 0.22{]} \\
\texttt{llm-b} & GPT & D5 & 0.00 & -0.27 & 0.27 {[}0.10, 0.46{]} \\
\texttt{llm-b} & GPT & D6 & -0.15 & -0.16 & 0.01 {[}-0.34, 0.33{]} \\
\texttt{llm-b} & GPT & D7 & -0.10 & -0.02 & -0.08 {[}-0.29, 0.05{]} \\
\texttt{llm-b} & GPT & D8 & -0.15 & 0.03 & -0.18 {[}-0.43, 0.05{]} \\
\texttt{llm-b} & GPT & D9 & 0.30 & -0.26 & 0.56 {[}0.18, 0.99{]} \\
\end{longtable}
}

Both LLM raters show self-preference on the total, with a CI above zero: \texttt{llm-a} by 1.11
points of 18 on Claude answers and \texttt{llm-b} by 1.50 points on GPT answers. Each scores the
other five families lower than the other two raters do, but scores its own vendor's family
about level with them. The gaps are of the same size as the mode effect on the total
(+1.37). For \texttt{llm-a} the preference sits in D2, D3 and D6. For \texttt{llm-b} it sits in D1, D2,
D3, D5 and D9. \texttt{llm-b}'s overall family guesses did not beat chance (Table 7), but whether
it recognised GPT answers in particular is not tested here.

Because the self-preference is as large as the mode effect, we ran one further check,
exploratory and not pre-registered (Appendix D): H4 with each LLM rater's own-vendor scores
replaced by the mean of the other two raters (\texttt{results/own-vendor-replaced/}). The verdict is
unchanged, model \textgreater{} mode: delta -0.506 {[}-0.646, -0.400{]}, omega squared 0.040 for mode and
0.546 for family, and reasoning minus fast 1.17 {[}0.58, 1.73{]}. Self-preference does not drive
the H4 result.

\section{Discussion}

\subsection{What the result means for evaluation practice}

On this task, which model answered mattered far more than how hard it was told to think.
Family explained 58\% of the variance in the total score after bias correction, mode 5\%,
and the interval on their difference lies wholly below zero (Section 4.2). Every dimension
points the same way, and the verdict survives rank transformation, each rater's own
scores, both fast-mode groups and the removal of each LLM rater's own-vendor scores. That
runs against the observation that motivated the study: LA-P1 saw one family's answers
change across modes and asked whether mode could rival the choice of vendor. In a
controlled grid it does not.

The result does not say mode is irrelevant. Reasoning scored 1.37 points higher than fast
out of 18, with an interval above zero. But that average hides most of the story. Three of
six families improved, one did not move and two got worse (Table 5). Qwen alone rose by
6.80 points; the other five averaged +0.28. The gain sits almost entirely in the families
whose fast mode turns reasoning off (+2.53), not in those where it only lowers it (+0.20).
And it is tied to length: once answer length is removed, the mode effect on the total
falls to 0.62 points and mode's share of the variance to 0.008 (Section 4.5). Reasoning
answers are longer, and longer answers score higher on this rubric. The design cannot say
whether the extra length carries the extra quality or merely accompanies it.

For anyone comparing models, the practical reading is twofold. First, the choice of model
should come first: no mode setting closed the gap between the strongest and weakest
families (GPT at 15.0 and 16.4, Gemini at 9.4 in both modes). Second, mode should
still be reported, because its effect is not one number. The family x mode interaction
(omega squared 0.184) is larger than the average mode effect (0.050), so ``reasoning helps''
is true for some models and false for others. A leaderboard that names the mode lets a
reader see that; one that does not hides a difference of up to 6.8 points for a single
model.

The LLM raters add a caution of their own. Each scored its own vendor's answers about 1.1
to 1.5 points higher, relative to the other raters, than it scored everyone else (Section
4.8). That is as large as the average mode effect. The H4 verdict does not depend on it,
but a study that compares modes or models with a single LLM judge from one of the vendors
being compared could easily mistake self-preference for an effect.

\subsection{Mode-sensitive and mode-stable dimensions}

No dimension was mode-sensitive in the pre-registered sense. Family dominated eight of
nine, and the ninth (D7, actionability) was inconclusive. The secondary hypotheses that
predicted a mode effect, H1 on hardware-constraint adherence and H2 on internal
consistency, both failed in direction, and all three secondary tests are exploratory
because their dimensions did not reach the reliability bar.

Within that pattern, the dimensions differ in how much mode moved them. D9, internal
consistency, had the largest reasoning gain (+0.43 of 2) and the only gain that survived
the length adjustment (+0.41). It also had the largest interaction (eta squared 0.226),
as Claude (0.4 to 1.4) and Qwen (0.0 to 1.8) rose while Grok fell (2.0 to 1.2). If any part of the answer
benefits from reasoning beyond writing more, it is keeping the deliverables consistent with
each other. D9's agreement was low (alpha 0.29), so this is a lead for future work, not a
finding. D1, the memory-budget check that H1 expected reasoning to improve, did not move at
all on average (0.00), even though it is the dimension where a careful calculation should
help most.

D3, tool factuality, was mode-stable, as H3 predicted. The competing account, that
reasoning models invent more tools, did not appear: reasoning minus fast on D3 was -0.10
with an interval that includes zero. This dimension carries a reliability warning of its
own. Its alpha was the lowest of all (0.25), and the human rater's ranking of the 12 cells
on D3 was unrelated to either LLM rater's (Section 4.7), even though all three
web-verified every tool. What counts as a real, correctly described tool is where the
raters disagreed most.

The two dimensions most closely tied to family, D2 (recency, eta squared 0.69) and D6
(citation quality, 0.65), are the ones that depend most on current information. The three
families whose apps searched the web on their own (Claude, GPT, Grok) scored 1.4 to 2.0 on
D2, and the families that ran without search (Gemini, DeepSeek) and Qwen scored 1.0 to 1.2.
This is consistent with live search driving part of the family effect, but search was not
varied within a family, so the study cannot separate it from the model itself (Section
5.3).

\subsection{Limitations}

\textbf{One task.} All 60 answers respond to one prompt. Replicates measure variability on that
task, not across design tasks. The task is long, open-ended and constrained by a stated
memory budget, so it may favour careful reasoning more than a typical request does.

\textbf{Six families, treated as fixed.} The conclusions are about these six models in these
settings during one 26-hour collection window, not about models in general. The models are
not matched on size or price, and some are not their vendor's strongest model.

\textbf{Mode is not one control.} ``Fast'' means reasoning switched off for GPT, Qwen and
DeepSeek, and the lowest available level for Claude, Gemini and Grok. Even an off switch
may not stop all reasoning on an open-ended task. The API check shows each switch changes
how much the model reasons, but it cannot show what a chat app runs internally. The
pre-registered split by fast-mode floor shows whether a result holds only for off versus
on.

\textbf{Chat apps, not the API.} Each vendor's app adds its own system prompt and may change
without notice. This matches how people use these models and how LA-P1 collected its
answers, but it trades away control. In particular, three apps (Claude, GPT and Grok)
searched the web by themselves with no way to turn it off, while two (Gemini and DeepSeek)
ran with search off. Differences between families may partly reflect live search.

\textbf{Model identity.} Chat apps change their models without asking. Two of six apps had
moved to a newer model by the time of collection (GPT-6 and Grok 4.7 instead of the
pre-registered GPT-5.6 Sol and Grok 4.6). Each still ran one model in both modes, so the
contrasts stay mode switches, but the API evidence for those two switches comes from the
planned models. DeepSeek's app shows no model name at all. This is the cost of testing the
apps people use rather than a fixed API snapshot, and the reason every capture is dated.

\textbf{Family is confounded with date and operator.} Families were captured one block at a
time, so family is partly confounded with time of day, and three blocks were sent by an AI
browser assistant and three by the author. Neither is confounded with mode, because fast
and reasoning runs were interleaved within each block.

\textbf{Blinding is partial.} Raters see recommended models and tools, which D2 to D4 need, and
an answer that favours its own vendor's models may reveal its family. Reasoning answers may
also read differently. We measure this as the rate at which raters guess family and mode
above chance, and the pre-registration sets what follows if they do.

\textbf{LLM raters.} Two of three raters are LLMs from vendors whose families are scored. We
report each one's scores for its own vendor's family separately, report agreement between
the human and each LLM, and use the median of three. A single human rater limits what the
study can say about human judgment. Both LLM raters did score their own vendor's family
higher than the other raters did (Section 4.8); removing those scores leaves H4 unchanged.
Two of three raters, the human and \texttt{llm-a}, guessed mode above chance, so mode blinding was
partial for them (Section 4.6).

\textbf{Answer length.} Reasoning answers are longer, and after removing log length the mode
effect on the total falls by more than half (Section 4.5). As pre-registered, we report
that the mode effect cannot be separated from answer length. We do not claim length causes
it, and we do not claim it does not.

\textbf{Small cells.} Five replicates per cell detect a strong mode effect reliably but miss a
moderate one about half the time, and percentile bootstrap intervals at this size can be
too narrow. Under equal effects, 15\% of simulated studies still show ``mode \textgreater{} model''. We
report permutation p-values beside every interval.

\textbf{Single author.} No second author checked the analysis. The plan was registered before
collection, the author did not score, and all data and code are released so anyone can
re-run it.

\section{Conclusion}

We crossed six model families with two reasoning modes on one demanding, constrained design
task, collected five answers per cell in the vendors' own chat apps, and had three blind
raters score all 60. On the pre-registered test, family explained far more of the variance
in answer quality than mode did (delta omega squared -0.527, 95\% CI {[}-0.666, -0.417{]}).
Reasoning helped on average, by 1.37 points of 18, but the gain came mostly from one family
and from switching reasoning off rather than turning it down, and it halved once answer
length was taken into account. On this task, choosing the model matters more than choosing
how hard it thinks. Mode should still be reported, because what it does differs from one
model to the next. The data, code and every departure from the plan are released.


\begin{acks}
We thank the human rater, who scored all 60 answers by hand and asked to remain anonymous.
[Funding statement: to be confirmed by the author.]
\end{acks}

\printbibliography

\appendix

% Generated by paper/latex_build.py from paper/*.md. Do not edit by hand.

\section{Mode registry}

What each mode was, how it was switched, and the evidence. Chat settings were used for the
60 scored answers. API settings were used once per cell to show the switch changes how much
the model reasons. The API check does not show what a chat app runs internally.

{\def\LTcaptype{none} % do not increment counter
\begin{longtable}[]{@{}
  >{\raggedright\arraybackslash}p{(\linewidth - 14\tabcolsep) * \real{0.1250}}
  >{\raggedright\arraybackslash}p{(\linewidth - 14\tabcolsep) * \real{0.1250}}
  >{\raggedright\arraybackslash}p{(\linewidth - 14\tabcolsep) * \real{0.1250}}
  >{\raggedright\arraybackslash}p{(\linewidth - 14\tabcolsep) * \real{0.1250}}
  >{\raggedright\arraybackslash}p{(\linewidth - 14\tabcolsep) * \real{0.1250}}
  >{\raggedright\arraybackslash}p{(\linewidth - 14\tabcolsep) * \real{0.1250}}
  >{\raggedright\arraybackslash}p{(\linewidth - 14\tabcolsep) * \real{0.1250}}
  >{\raggedright\arraybackslash}p{(\linewidth - 14\tabcolsep) * \real{0.1250}}@{}}
\toprule\noalign{}
\begin{minipage}[b]{\linewidth}\raggedright
Family
\end{minipage} & \begin{minipage}[b]{\linewidth}\raggedright
Chat app, fast
\end{minipage} & \begin{minipage}[b]{\linewidth}\raggedright
Chat app, reasoning
\end{minipage} & \begin{minipage}[b]{\linewidth}\raggedright
API model
\end{minipage} & \begin{minipage}[b]{\linewidth}\raggedright
API fast
\end{minipage} & \begin{minipage}[b]{\linewidth}\raggedright
API reasoning
\end{minipage} & \begin{minipage}[b]{\linewidth}\raggedright
Docs read
\end{minipage} & \begin{minipage}[b]{\linewidth}\raggedright
Reasoning tokens (fast vs reasoning)
\end{minipage} \\
\midrule\noalign{}
\endhead
\bottomrule\noalign{}
\endlastfoot
Claude & claude.ai, Sonnet 5.5, Effort Low & claude.ai, Sonnet 5.5, Effort High & \texttt{claude-sonnet-5-5} & adaptive thinking, effort low & adaptive thinking, effort high & Anthropic API reference, 2026-10-06 & 0 vs 14,864 \\
GPT & chatgpt.com (paid), GPT-6, Thinking effort Instant & same, Thinking effort High & \texttt{gpt-5.6-sol} & \texttt{reasoning\_effort:\ none} & \texttt{reasoning\_effort:\ high} & OpenAI model page, 2026-10-06 & 0 vs 6,000 \\
Gemini & aistudio.\allowbreak{}google.com, Gemini 3.8 Flash, Thinking level Low, Search grounding off & same, Thinking level High & \texttt{gemini-3.8-flash} & \texttt{reasoning\_effort:\ low} & \texttt{reasoning\_effort:\ high} & Gemini thinking docs, 2026-10-13 & 0 vs 4,911 \\
Qwen & chat.qwen.ai, Qwen3.8-Max, Fast & same, Thinking & \texttt{qwen3.8-max} & \texttt{enable\_thinking:\ false} & \texttt{enable\_thinking:\ true} & Alibaba Model Studio docs, 2026-10-06 & 0 vs 23,805 \\
DeepSeek & chat.\allowbreak{}deepseek.com, DeepThink off, Search off & same, DeepThink on & \texttt{deepseek-v4.1-flash} (Alibaba-hosted) & \texttt{enable\_thinking:\ false} & \texttt{enable\_thinking:\ true} & Alibaba Model Studio docs, 2026-10-13 & 0 vs 1,367 \\
Grok & grok.com, Grok 4.7, Fast & grok.com, Grok 4.7, Expert & \texttt{grok-4.6} & \texttt{reasoning\_effort:\ low} & \texttt{reasoning\_effort:\ xhigh} & xAI reasoning docs, 2026-10-05 & 131 vs 4,627 \\
\end{longtable}
}

Notes on the evidence:

\begin{itemize}
\tightlist
\item
  \textbf{Fast floor.} GPT, Qwen and DeepSeek can switch reasoning fully off. Claude Sonnet 5.5
  rejects \texttt{thinking:\ disabled}, Gemini 3.8 Flash has no level below low, and Grok 4.7
  cannot disable reasoning. For those three, ``fast'' is the lowest available setting. Claude
  and Gemini returned zero reasoning tokens at that setting on this one call, but we keep
  them in the lowest-available group because their documentation says reasoning cannot be
  turned off.
\item
  \textbf{Limits of the API calls.} The GPT reasoning call spent its whole 6,000-token output
  limit, set by the account's rate limit, on reasoning and returned no visible answer. The
  Claude reasoning call stopped at its 32,000-token limit. The DeepSeek and Qwen calls ran
  on Alibaba Model Studio, a third-party host. None of this affects the evidence that the
  switch works.
\item
  \textbf{Model identity in the app.} Claude, GPT, Gemini and Qwen show the model name in the
  model picker. The author confirmed that grok.com ran Grok 4.7 in both modes. DeepSeek
  shows no model name, so its model is the author's observation.
\item
  \textbf{Planned and actual models.} The pre-registration listed GPT-5.6 Sol and Grok 4.6. GPT-6
  reached ChatGPT on 7 October 2026 and Grok 4.7 reached grok.com in early October, and the
  apps served those during collection. The API checks for these two families ran on the
  planned API models, so they show the switch on GPT-5.6 Sol and Grok 4.6 only. GPT-6 Sol
  accepts \texttt{reasoning.effort:\ none} in OpenAI's API, and Grok 4.7 cannot disable reasoning,
  so neither family changes group in robustness check 4. No release note for any of the six
  models was found inside the collection window. Sources are in the mode registry.
\item
  \textbf{Superseded checks.} Earlier API checks for Gemini (on 3.6 Flash and before an adapter
  fix) and DeepSeek (on V4 Pro) are kept in \texttt{data/captures/superseded/} with the reason for
  each.
\end{itemize}

\section{Per-cell scores}

Consensus means per family and mode, five answers per cell (\texttt{results/cell\_means.csv}).
Each dimension is scored 0 to 2, the total 0 to 18. No cell has an answer scored 0 on all
nine dimensions.

{\def\LTcaptype{none} % do not increment counter
\begin{longtable}[]{@{}
  >{\raggedright\arraybackslash}p{(\linewidth - 22\tabcolsep) * \real{0.0833}}
  >{\raggedright\arraybackslash}p{(\linewidth - 22\tabcolsep) * \real{0.0833}}
  >{\raggedright\arraybackslash}p{(\linewidth - 22\tabcolsep) * \real{0.0833}}
  >{\raggedright\arraybackslash}p{(\linewidth - 22\tabcolsep) * \real{0.0833}}
  >{\raggedright\arraybackslash}p{(\linewidth - 22\tabcolsep) * \real{0.0833}}
  >{\raggedright\arraybackslash}p{(\linewidth - 22\tabcolsep) * \real{0.0833}}
  >{\raggedright\arraybackslash}p{(\linewidth - 22\tabcolsep) * \real{0.0833}}
  >{\raggedright\arraybackslash}p{(\linewidth - 22\tabcolsep) * \real{0.0833}}
  >{\raggedright\arraybackslash}p{(\linewidth - 22\tabcolsep) * \real{0.0833}}
  >{\raggedright\arraybackslash}p{(\linewidth - 22\tabcolsep) * \real{0.0833}}
  >{\raggedright\arraybackslash}p{(\linewidth - 22\tabcolsep) * \real{0.0833}}
  >{\raggedright\arraybackslash}p{(\linewidth - 22\tabcolsep) * \real{0.0833}}@{}}
\toprule\noalign{}
\begin{minipage}[b]{\linewidth}\raggedright
Family
\end{minipage} & \begin{minipage}[b]{\linewidth}\raggedright
Mode
\end{minipage} & \begin{minipage}[b]{\linewidth}\raggedright
D1
\end{minipage} & \begin{minipage}[b]{\linewidth}\raggedright
D2
\end{minipage} & \begin{minipage}[b]{\linewidth}\raggedright
D3
\end{minipage} & \begin{minipage}[b]{\linewidth}\raggedright
D4
\end{minipage} & \begin{minipage}[b]{\linewidth}\raggedright
D5
\end{minipage} & \begin{minipage}[b]{\linewidth}\raggedright
D6
\end{minipage} & \begin{minipage}[b]{\linewidth}\raggedright
D7
\end{minipage} & \begin{minipage}[b]{\linewidth}\raggedright
D8
\end{minipage} & \begin{minipage}[b]{\linewidth}\raggedright
D9
\end{minipage} & \begin{minipage}[b]{\linewidth}\raggedright
Total
\end{minipage} \\
\midrule\noalign{}
\endhead
\bottomrule\noalign{}
\endlastfoot
Claude & fast & 1.0 & 2.0 & 2.0 & 2.0 & 1.6 & 1.6 & 1.2 & 2.0 & 0.4 & 13.8 \\
Claude & reasoning & 1.2 & 2.0 & 2.0 & 1.8 & 2.0 & 1.6 & 1.6 & 2.0 & 1.4 & 15.6 \\
DeepSeek & fast & 1.4 & 1.0 & 1.8 & 1.8 & 1.2 & 0.8 & 1.2 & 1.8 & 0.8 & 11.8 \\
DeepSeek & reasoning & 1.2 & 1.0 & 1.4 & 2.0 & 1.4 & 0.6 & 0.8 & 1.8 & 1.0 & 11.2 \\
Gemini & fast & 1.4 & 1.0 & 1.4 & 1.8 & 0.0 & 1.0 & 1.0 & 1.6 & 0.2 & 9.4 \\
Gemini & reasoning & 1.4 & 1.0 & 1.0 & 1.8 & 0.0 & 1.0 & 1.0 & 1.8 & 0.4 & 9.4 \\
GPT & fast & 1.8 & 1.4 & 2.0 & 2.0 & 2.0 & 1.8 & 1.0 & 1.4 & 1.6 & 15.0 \\
GPT & reasoning & 1.2 & 2.0 & 1.8 & 2.0 & 2.0 & 2.0 & 1.8 & 1.8 & 1.8 & 16.4 \\
Grok & fast & 1.0 & 2.0 & 2.0 & 2.0 & 1.6 & 1.0 & 1.0 & 1.0 & 2.0 & 13.6 \\
Grok & reasoning & 1.0 & 1.8 & 2.0 & 2.0 & 1.4 & 1.0 & 1.0 & 1.0 & 1.2 & 12.4 \\
Qwen & fast & 0.6 & 1.2 & 1.0 & 0.6 & 0.4 & 1.0 & 1.0 & 1.0 & 0.0 & 6.8 \\
Qwen & reasoning & 1.2 & 1.2 & 1.4 & 1.6 & 2.0 & 1.0 & 1.4 & 2.0 & 1.8 & 13.6 \\
\end{longtable}
}

Reasoning minus fast per family and dimension is in \texttt{results/family\_mode\_effects.csv}.
The full decomposition is Table 4. The random-effects variance components, exploratory as
pre-registered, are below.

{\def\LTcaptype{none} % do not increment counter
\begin{longtable}[]{@{}lllll@{}}
\toprule\noalign{}
Dim & Family & Mode & Family x mode & Residual \\
\midrule\noalign{}
\endhead
\bottomrule\noalign{}
\endlastfoot
D1 & 0.016 & 0.000 & 0.043 & 0.183 \\
D2 & 0.188 & 0.000 & 0.022 & 0.075 \\
D3 & 0.120 & 0.000 & 0.024 & 0.108 \\
D4 & 0.072 & 0.000 & 0.068 & 0.117 \\
D5 & 0.389 & 0.020 & 0.168 & 0.225 \\
D6 & 0.200 & 0.000 & 0.000 & 0.108 \\
D7 & 0.000 & 0.005 & 0.065 & 0.117 \\
D8 & 0.077 & 0.023 & 0.054 & 0.117 \\
D9 & 0.077 & 0.029 & 0.317 & 0.350 \\
total & 4.136 & 0.233 & 3.872 & 1.658 \\
\end{longtable}
}

\section{Pilot data and why it is excluded}

The study grew out of three DeepSeek answers collected for LA-P1 \cite{aryal2026lap1} on
2026-08-31, before this study was designed: one per mode (deep reasoning, fast, and fast
with the DeepThink toggle), scored with the v1 rubric and adjudicated by hand. Their scores
belong to LA-P1 and are not reproduced here.

The pilot is excluded from every analysis for three reasons: it predates the
pre-registration, it used the v1 rubric without the v2 anchors, and it has one answer per
mode, so it cannot separate a mode effect from run-to-run noise. As LA-P1 reports, its
adjudicated totals did not show a reasoning advantage; the pre-registration records this so
that the present result cannot be framed as confirming the pilot. LA-P1's early
fabrication audit also flagged real tools released after the auditor's training cutoff as
invented, which is why D3 and D4 here require web verification.

\section{Deviations from the pre-registration}

The pre-registration was frozen on 2026-10-13 and registered on OSF the same day. Every
later change is listed with its date and reason in section 10 of
\texttt{docs/preregistration.md}, which is reproduced in full in the released materials. In
summary:

{\def\LTcaptype{none} % do not increment counter
\begin{longtable}[]{@{}
  >{\raggedright\arraybackslash}p{(\linewidth - 4\tabcolsep) * \real{0.3333}}
  >{\raggedright\arraybackslash}p{(\linewidth - 4\tabcolsep) * \real{0.3333}}
  >{\raggedright\arraybackslash}p{(\linewidth - 4\tabcolsep) * \real{0.3333}}@{}}
\toprule\noalign{}
\begin{minipage}[b]{\linewidth}\raggedright
Area
\end{minipage} & \begin{minipage}[b]{\linewidth}\raggedright
Changes
\end{minipage} & \begin{minipage}[b]{\linewidth}\raggedright
Effect on the analysis
\end{minipage} \\
\midrule\noalign{}
\endhead
\bottomrule\noalign{}
\endlastfoot
Operators and send times & Three blocks sent by an AI browser assistant at the author's request, three by the author. Send times recorded as exact times or as bounded intervals where chats ran side by side & None on scores. Disclosed for transparency \\
Replacements & 6 failed attempts replaced: 3 Claude reasoning runs stopped by a usage limit (treated as the app being down), 1 DeepSeek cut-off, 1 DeepSeek app error, 1 Qwen run sent in the wrong mode (treated like a wrong model). One extra Qwen chat kept, unscored & 60 valid answers from 66 attempts. No family dropped \\
App behaviour & Claude, GPT and Grok searched the web by themselves, with no switch to stop it. Two DeepSeek fast answers were finished with the app's own ``Continue'' button & Family differences may include live search. Stated as a limitation \\
Models that ran & GPT ran GPT-6 and Grok ran Grok 4.7, not the pre-registered GPT-5.6 Sol and Grok 4.6, confirmed by the author. Each ran one model in both modes. The GPT captures' \texttt{model\_returned} field is corrected in the Deviations table, not in the write-once captures. Rater \texttt{llm-b} is GPT-6 too & API checks for those two families cover the planned models only. No change of analysis group \\
Model self-reports & DeepSeek and Gemini answers named Claude models, which is wrong. GPT (``GPT-6'') and Grok (``Grok 4.7'') answers named themselves correctly & DeepSeek contrast described as a vendor chat setting \\
Reasoning traces & Kept where the app showed them. One GPT trace not captured & None. Raters never see traces \\
Blinding, before scoring & Four fixes: chat-app clutter removed, bare ``Model:'' label lines removed, removed header lines leave no marker, and self-disclosures naming the answer's maker removed & Same rule for every packet. Key unchanged \\
Raters & Two LLM raters named after the freeze, as the plan requires. One answer per session from a fixed kit, instructions hashed before scoring. The first human rater's returns were made with AI tools and are not used; Rater H2 scored by hand and was recorded after scoring. 18 \texttt{llm-a} sessions worked the memory check by hand; the author re-ran every input and all results match & None on the analysis plan. Rater H2's scores are the human scores \\
Analysis, after scoring & The length check (robustness check 2) first judged shrinking on delta instead of the mode effect; corrected, and the mode effect is reported as not separable from length. One exploratory check added, not pre-registered: H4 with own-vendor LLM scores replaced & No pre-registered number changed. H4 verdict unchanged in the added check \\
\end{longtable}
}

\subsection{Pre-scoring blinding audit}

\texttt{results/blinding\_audit.csv}, from \texttt{modegrid\ blinding-audit} on the final packets, before any
scoring. Each row counts the packets in each cell that contain a pattern.

{\def\LTcaptype{none} % do not increment counter
\begin{longtable}[]{@{}
  >{\raggedright\arraybackslash}p{(\linewidth - 6\tabcolsep) * \real{0.2500}}
  >{\raggedright\arraybackslash}p{(\linewidth - 6\tabcolsep) * \real{0.2500}}
  >{\raggedright\arraybackslash}p{(\linewidth - 6\tabcolsep) * \real{0.2500}}
  >{\raggedright\arraybackslash}p{(\linewidth - 6\tabcolsep) * \real{0.2500}}@{}}
\toprule\noalign{}
\begin{minipage}[b]{\linewidth}\raggedright
Pattern
\end{minipage} & \begin{minipage}[b]{\linewidth}\raggedright
Kind
\end{minipage} & \begin{minipage}[b]{\linewidth}\raggedright
Packets
\end{minipage} & \begin{minipage}[b]{\linewidth}\raggedright
Where
\end{minipage} \\
\midrule\noalign{}
\endhead
\bottomrule\noalign{}
\endlastfoot
Masked self-reference, self-disclosure, app tracking parameters, app citation markup, Grok citation tags, app images, tool file links & leak & 0 & none \\
Tool-progress text (``I'll start by verifying \ldots{}'') & style & 8 & Claude, 4 fast and 4 reasoning \\
Note that the real date differs from the prompt's & style & 8 & Claude, 4 fast and 4 reasoning \\
``As of my cutoff'' & style & 8 & DeepSeek, 4 fast and 4 reasoning \\
LaTeX arrows in prose & style & 7 & Gemini, 4 fast and 3 reasoning \\
Mentions a downloadable archive & style & 1 & GPT, reasoning \\
\end{longtable}
}

Style habits are the model's own text and stay, as the pre-registration requires. Each one
is confined to a single family but balanced across modes, so it can reveal family but not
mode. Family guesses above chance are expected and do not change any hypothesis's status
(pre-registration section 4).

The audit had no probe for the author's own name. A check after scoring found it, or the
author's home city, in three packets, all Claude reasoning answers (Section 4.6 and the
Deviations table).

\section{Data and code availability}

Code (MIT) and data (CC-BY-4.0) are in tag \texttt{v1.0.0} of
\url{https://github.com/roshanaryal1/reasoning-mode-grid}, archived at
\url{https://doi.org/10.5281/zenodo.23674817}. The release holds the
configuration, the randomised manifest, all 72 captures (60 scored, 12 API checks) and the
superseded attempts, the blinding rules (packets are rebuilt from the captures and the
committed key), the three rater sheets, the consensus scores and every analysis output.
The prompt, rubric and memory checker are vendored from LA-P1 at commit \texttt{d5e39e7} and
pinned by SHA-256. All randomness is seeded.


\section{Reproducibility Checklist for JAIR}

Select the answers that apply to your research -- one per item.

\subsection*{All articles:}

\begin{enumerate}
    \item All claims investigated in this work are clearly stated.
    \textbf{yes}
    \item Clear explanations are given how the work reported substantiates the claims.
    \textbf{yes}
    \item Limitations or technical assumptions are stated clearly and explicitly.
    \textbf{yes}
    \item Conceptual outlines and/or pseudo-code descriptions of the AI methods introduced in this work are provided, and important implementation details are discussed.
    \textbf{yes}
    \item
    Motivation is provided for all design choices, including algorithms, implementation choices, parameters, data sets and experimental protocols beyond metrics.
    \textbf{yes}
\end{enumerate}

\subsection*{Articles containing theoretical contributions:}
Does this paper make theoretical contributions?
\textbf{no}

\subsection*{Articles reporting on computational experiments:}
Does this paper include computational experiments? \textbf{yes}

\begin{enumerate}
    \item
    All source code required for conducting experiments is included in an online appendix
    or will be made publicly available upon publication of the paper.
    The online appendix follows best practices for source code readability and documentation as well as for long-term accessibility.
    \textbf{yes}
    \item The source code comes with a license that
    allows free usage for reproducibility purposes.
    \textbf{yes}
    \item The source code comes with a license that
    allows free usage for research purposes in general.
    \textbf{yes}
    \item
    Raw, unaggregated data from all experiments is included in an online appendix
    or will be made publicly available upon publication of the paper.
    The online appendix follows best practices for long-term accessibility.
    \textbf{yes}
    \item The unaggregated data comes with a license that
    allows free usage for reproducibility purposes.
    \textbf{yes}
    \item The unaggregated data comes with a license that
    allows free usage for research purposes in general.
    \textbf{yes}
    \item If an algorithm depends on randomness, then the method used for generating random numbers and for setting seeds is described in a way sufficient to allow replication of results.
    \textbf{yes}
    \item The execution environment for experiments, the computing infrastructure (hardware and software) used for running them, is described, including GPU/CPU makes and models; amount of memory (cache and RAM); make and version of operating system; names and versions of relevant software libraries and frameworks.
    \textbf{partially}
    \item
    The evaluation metrics used in experiments are clearly explained and their choice is explicitly motivated.
    \textbf{yes}
    \item
    The number of algorithm runs used to compute each result is reported.
    \textbf{yes}
    \item
    Reported results have not been ``cherry-picked'' by silently ignoring unsuccessful or unsatisfactory experiments.
    \textbf{yes}
    \item
    Analysis of results goes beyond single-dimensional summaries of performance (e.g., average, median) to include measures of variation, confidence, or other distributional information.
    \textbf{yes}
    \item
    All (hyper-) parameter settings for
    the algorithms/methods used in experiments have been reported, along with the rationale or method for determining them.
    \textbf{yes}
    \item
    The number and range of (hyper-) parameter settings explored prior to conducting final experiments have been indicated, along with the effort spent on (hyper-) parameter optimisation.
    \textbf{NA}
    \item
    Appropriately chosen statistical hypothesis tests are used to establish statistical significance
    in the presence of noise effects.
    \textbf{yes}
\end{enumerate}

\subsection*{Articles using data sets:}
Does this work rely on one or more data sets (possibly obtained from a benchmark generator or similar software artifact)?
\textbf{yes}

\begin{enumerate}
    \item
    All newly introduced data sets
    are included in an online appendix
    or will be made publicly available upon publication of the paper.
    The online appendix follows best practices for long-term accessibility with a license
    that allows free usage for research purposes.
    \textbf{yes}
    \item The newly introduced data set comes with a license that
    allows free usage for reproducibility purposes.
    \textbf{yes}
    \item The newly introduced data set comes with a license that
    allows free usage for research purposes in general.
    \textbf{yes}
    \item All data sets drawn from the literature or other public sources (potentially including authors' own previously published work) are accompanied by appropriate citations.
    \textbf{yes}
    \item All data sets drawn from the existing literature (potentially including authors' own previously published work) are publicly available. \textbf{yes}
    \item All new data sets and data sets that are not publicly available are described in detail, including relevant statistics, the data collection process and annotation process if relevant.
    \textbf{yes}
    \item
    All methods used for preprocessing, augmenting, batching or splitting data sets (e.g., in the context of hold-out or cross-validation)
    are described in detail. \textbf{yes}
\end{enumerate}

\subsection*{Explanations on any of the answers above (optional):}

Code (MIT licence) and data (CC-BY-4.0) are in a repository that will be made public, with
an archived release, on publication. The answers were produced inside vendors' chat apps,
whose hardware and serving software cannot be observed; the paper reports each app, model,
mode setting and capture time instead. The analysis runs in Python 3.10 or later; minimum
dependency versions are in \texttt{pyproject.toml}, and every random step uses a fixed seed.
No hyper-parameters were tuned: the mode settings and the analysis were fixed in the
pre-registration before data collection, and every later change is listed in the Deviations
appendix.

\end{document}
\endinput
